\section{训练目标其实是一个去噪回归}
\label{sec:14-Deep-Learning/08-Diffusion-Models/02}

你打开那份代码的训练循环，准备找变分下界。论文里那一页半的推导有 KL 项、有对数行列式、有一串下标滚来滚去的求和，你以为代码会长得同样吓人。实际看到的是五行：从 \(\{1,\dots,T\}\) 里均匀抽一个 \(t\)，抽一个标准高斯 \(\varepsilon\)，按上一节的闭式解合成 \(x_t\)，喂进网络，最后一行是对网络输出和 \(\varepsilon\) 求均方误差。整个仓库里没有第二个损失函数。

这中间发生的事情值得完整走一遍，不是为了欣赏代数，而是因为省掉的那几步里有一步\emph{不是}恒等变形。分不清哪些步骤是推导、哪些步骤是选择，你就说不清这个模型到底在优化什么。

\subsection{下界的形状和 VAE 是同一个}

把 \(x_1,\dots,x_T\) 整体看成隐变量，\(x_0\) 是观测，前向过程 \(q(x_{1:T}\given x_0)\) 就是一个\emph{不含参数}的推断分布。这个结构和 \hyperref[sec:14-Deep-Learning/06-Variational-Autoencoders/01]{``VAE 把不可积后验变成可训练下界''}里的完全一致，区别只在于那里的 \(q_\phi\) 要学，这里的 \(q\) 是写死的。同一个 Jensen 不等式给出同一个界：

\[
-\log p_\theta(x_0) \le \E_q\left[-\log\frac{p_\theta(x_{0:T})}{q(x_{1:T}\given x_0)}\right] \eqdef L .
\]

把分子分母按马尔可夫结构逐项配对、再对 \(q\) 取期望，\(L\) 可以重排成一串条件项：

\[
L = \underbrace{\KL\bigl(q(x_T\given x_0)\,\big\|\,p(x_T)\bigr)}_{L_T}
+ \sum_{t=2}^{T}\underbrace{\KL\bigl(q(x_{t-1}\given x_t,x_0)\,\big\|\,p_\theta(x_{t-1}\given x_t)\bigr)}_{L_{t-1}}
\underbrace{-\;\log p_\theta(x_0\given x_1)}_{L_0}.
\]

\(L_T\) 里没有 \(\theta\)：前向固定，先验也固定，它是一个常数，调度设计得当时接近零。\(L_0\) 是最后一步把连续量落回像素取值的解码项，实现上用一个离散化的高斯，占的比重很小。真正扛住整个训练的是中间那一串 \(L_{t-1}\)。

这个拆法把``一次难的分布匹配''换成了``\(T\) 次容易的分布匹配''。原问题是让 \(p_\theta\) 在整个数据空间上对齐一个极复杂的分布；拆完之后每一项只要求模型在\emph{已知 \(x_t\)} 的前提下匹配一小步的反向条件分布，而这一小步的跨度可以做得任意小。步子小到什么程度，那个条件分布就简单到什么程度。

\subsection{每一项都是两个高斯之间的 KL}

\(L_{t-1}\) 之所以有闭式，靠的是一个不那么显然的事实：虽然反向条件分布 \(q(x_{t-1}\given x_t)\) 一般不是高斯（要对整个数据分布积分才能得到它），但\emph{额外给定 \(x_0\)} 之后它就是高斯，而且均值方差都能写出来：

\[
q(x_{t-1}\given x_t,x_0) = \mathcal N\bigl(\tilde\mu_t(x_t,x_0),\;\tilde\beta_t I\bigr),
\qquad
\tilde\beta_t = \frac{1-\bar\alpha_{t-1}}{1-\bar\alpha_t}\,\beta_t,
\]

\[
\tilde\mu_t(x_t,x_0)
= \frac{\sqrt{\bar\alpha_{t-1}}\,\beta_t}{1-\bar\alpha_t}\,x_0
+ \frac{\sqrt{\alpha_t}\,(1-\bar\alpha_{t-1})}{1-\bar\alpha_t}\,x_t .
\]

证明骨架是一次 Bayes 配方：\(q(x_{t-1}\given x_t,x_0)\propto q(x_t\given x_{t-1})\,q(x_{t-1}\given x_0)\)，两个因子在指数上都是 \(x_{t-1}\) 的二次式，配平方即得。条件对账在于第二个因子——它需要 \(x_0\) 才写得出来，去掉 \(x_0\) 就得对数据分布积分，配方立刻失效。变分下界的可算性完全押在这一条上：训练时 \(x_0\) 是手上的那张图片，所以能算；采样时没有 \(x_0\)，所以要靠网络。

把 \(p_\theta(x_{t-1}\given x_t)\) 也取成高斯 \(\mathcal N(\mu_\theta(x_t,t),\sigma_t^2 I)\)，并把方差固定成常数（\(\sigma_t^2\) 取 \(\beta_t\) 或 \(\tilde\beta_t\) 都可以，实践上差别不大），两个同协方差高斯之间的 KL 就塌成一个平方距离：

\[
\KL\bigl(\mathcal N(\tilde\mu,\sigma^2 I)\,\big\|\,\mathcal N(\mu_\theta,\sigma^2 I)\bigr)
= \frac{1}{2\sigma^2}\,\norm{\tilde\mu-\mu_\theta}^2 .
\]

散度消失了，剩下均方误差。这是整条推导链上最关键的一次结构简化，代价是把 \(p_\theta\) 的每一步限死成各向同性高斯——步长足够小时这个限制不伤害表达力，因为小步的真实反向分布本来就接近高斯。

\subsection{重参数化：把``预测均值''改写成``预测噪声''}

现在只差把 \(\tilde\mu_t\) 写成模型看得见的量的函数。\(\tilde\mu_t\) 里含 \(x_0\)，而模型输入只有 \(x_t\)。用上一节那个闭式解 \(x_t=\sqrt{\bar\alpha_t}x_0+\sqrt{1-\bar\alpha_t}\,\varepsilon\) 反解出 \(x_0=(x_t-\sqrt{1-\bar\alpha_t}\,\varepsilon)/\sqrt{\bar\alpha_t}\)，代回去合并同类项：

\[
\begin{aligned}
\tilde\mu_t
&= \frac{\sqrt{\bar\alpha_{t-1}}\,\beta_t}{1-\bar\alpha_t}\cdot\frac{x_t-\sqrt{1-\bar\alpha_t}\,\varepsilon}{\sqrt{\bar\alpha_t}}
 + \frac{\sqrt{\alpha_t}\,(1-\bar\alpha_{t-1})}{1-\bar\alpha_t}\,x_t\\[2pt]
&= \frac{1}{\sqrt{\alpha_t}}\left(x_t-\frac{\beta_t}{\sqrt{1-\bar\alpha_t}}\,\varepsilon\right).
\end{aligned}
\]

这个形式很有说服力：反向一步的目标均值，等于把当前样本按 \(1/\sqrt{\alpha_t}\) 放大一点，再减掉里面那份噪声的一个固定倍数。要算出它，唯一缺的就是 \(\varepsilon\)。于是让网络去补这个缺口，令

\[
\mu_\theta(x_t,t)=\frac{1}{\sqrt{\alpha_t}}\left(x_t-\frac{\beta_t}{\sqrt{1-\bar\alpha_t}}\,\varepsilon_\theta(x_t,t)\right),
\]

两者相减时 \(x_t\) 项抵消，只剩系数乘以噪声之差：

\[
L_{t-1} = \frac{\beta_t^{2}}{2\sigma_t^{2}\,\alpha_t\,(1-\bar\alpha_t)}\,
\norm{\varepsilon-\varepsilon_\theta(x_t,t)}^{2}
\eqdef w_t\,\norm{\varepsilon-\varepsilon_\theta(x_t,t)}^{2}.
\]

到这一步为止，每一步都是恒等变形或者明确声明过的建模选择。变分下界已经变成了一个逐时刻加权的最小二乘。

\subsection{丢掉权重的那一步不是推导}

实践中用的不是上式，而是把 \(w_t\) 全部换成 \(1\)：

\[
L_{\mathrm{simple}}(\theta)=
\E_{t,\;x_0,\;\varepsilon}\,
\norm{\varepsilon-\varepsilon_\theta(x_t,t)}^{2},
\qquad
t\sim\Uniform\{1,\dots,T\},\;
x_0\sim p_{\mathrm{data}},\;
\varepsilon\sim\mathcal N(0,I).
\]

训练一个扩散模型，做的就是这件事：随机取一个噪声水平，把干净数据按那个水平弄脏，让网络回归它加进去的那份噪声。没有对抗博弈，没有可逆性约束，没有配分函数。生成建模里最难的那一类问题，在这里被写成了一个监督回归。

而丢掉 \(w_t\) 这件事必须说清楚：\textbf{它不是恒等变形，也不是放缩，去掉权重之后的目标不再是似然的下界。}\(w_t\) 在小 \(t\)（低噪声）处非常大——那里 \(1-\bar\alpha_t\) 接近零，分母把权重推上去——所以严格的下界会把绝大部分训练资源压在几乎没被破坏的样本上，让网络去抠那些人眼根本分辨不出、对生成质量贡献极小的最后一点像素误差。均匀加权把资源挪回中高噪声区段，也就是上一节说的``学得最多''的那一段。

这个替换带来的是一次有方向的交易，两侧都能观测到：按 \(L_{\mathrm{simple}}\) 训练出来的模型，样本质量指标明显更好，而似然指标（每维比特数）明显更差。后一半不是缺陷，是定义的直接后果——你优化的已经不是似然了，似然变差理所当然。选择 \(L_{\mathrm{simple}}\) 是一条经验规律，靠的是在图像基准上反复比较，不是任何定理保证的；换个数据模态、换个评价指标，最优的加权方式完全可能不同，后来确实有一批工作专门去重新设计 \(w_t\)。

\begin{remark}
把这件事放在书里反复出现的那条分界线上看：``下界''是一个可以证明的陈述，``重新加权的下界效果更好''是一个经验陈述。前者告诉你优化它一定不会让似然变坏，后者只告诉你在某些数据集上人们量到了更好的分数。混淆两者的代价是，当你的数据不属于那些数据集时，你会以为自己有理论保障，其实没有。
\end{remark}

\subsection{网络学到的东西就是得分}

上一节结尾说过，反向走这条路需要一个额外的量，那个量叫得分。现在可以把它和 \(\varepsilon_\theta\) 对上。

对固定的 \(x_0\)，\(q(x_t\given x_0)\) 是高斯，对数密度关于 \(x_t\) 的梯度直接可算：

\[
\nabla_{x_t}\log q(x_t\given x_0)
= -\frac{x_t-\sqrt{\bar\alpha_t}\,x_0}{1-\bar\alpha_t}
= -\frac{\varepsilon}{\sqrt{1-\bar\alpha_t}} .
\]

边缘密度 \(p_t(x)=\int q(x\given x_0)p_{\mathrm{data}}(x_0)\,dx_0\) 的得分则是这些条件得分的加权平均，权重恰好是后验 \(p(x_0\given x_t=x)\)，于是

\[
\nabla_x\log p_t(x) = -\frac{\E[\varepsilon\given x_t=x]}{\sqrt{1-\bar\alpha_t}} .
\]

另一边，平方损失的最优预测是条件期望——\hyperref[sec:05-Probability/04-Expectation-and-Moments/04]{``条件期望是利用信息的最佳预测''}证过这件事，且与预测器的具体形式无关。所以在网络容量足够、训练充分的理想情形下，\(\varepsilon_\theta(x,t)\) 收敛到的正是 \(\E[\varepsilon\given x_t=x]\)。两式一比：

\[
\nabla_x\log p_t(x)\;\approx\;-\frac{\varepsilon_\theta(x,t)}{\sqrt{1-\bar\alpha_t}} .
\]

去噪和估计得分是同一件事，只差一个已知的缩放。\hyperref[sec:06-Random-Processes/06-Diffusion-and-SDE/05]{``得分匹配：怎样从样本学习得分''}从另一头走到了同一个地方：那里的目标是直接拟合 \(\nabla\log p\)，代价是原始形式含有一个难算的散度项，去噪得分匹配把它换成了一个可以采样估计的形式——换出来的东西就是上面这个回归。

同样的关系还有一个更直观的说法，即 Tweedie 公式：对高斯加噪的观测，最优的去噪结果可以用观测点处的得分表示出来，

\[
\E[x_0\given x_t=x]=\frac{1}{\sqrt{\bar\alpha_t}}\Bigl(x+(1-\bar\alpha_t)\,\nabla_x\log p_t(x)\Bigr).
\]

读法是：从脏样本出发，沿着对数密度上升最快的方向挪一段，挪的距离正比于噪声方差，落点就是最优的干净估计。这句话把下一节的采样算法提前解释掉了一半——反向过程之所以能从噪声走回数据，是因为每一步都在沿着得分往密度高的地方走，而得分正是训练好的网络（缩放后）直接输出的东西。噪声大的时候密度被抹得很平，得分场很平缓，只能给出粗略的方向；噪声小的时候密度重新变得尖锐，得分场把样本精确地拉向数据流形。

\subsection{一个网络要同时干很多件不同的活}

\(t\) 是网络的输入，不是每个 \(t\) 各训一个网络。这个实现细节的后果比它看起来大：\(L_{\mathrm{simple}}\) 里对 \(t\) 取期望，意味着同一组参数要在\emph{所有}噪声水平上同时把回归做好，而不同噪声水平上的``最优去噪''是性质完全不同的任务。

\begin{center}
\begin{tikzpicture}[font=\small,>=stealth]
% ===== 三个噪声水平的一维信号 =====
% --- 低噪声 ---
\begin{scope}[shift={(0,2.3)}]
  \draw[black!12] (0,0) rectangle (2.8,1.2);
  \draw[black!30,thin,smooth] plot coordinates
    {(0,0.70)(0.2,0.82)(0.4,0.94)(0.6,1.00)(0.8,0.94)(1.0,0.76)(1.2,0.54)
     (1.4,0.36)(1.6,0.26)(1.8,0.30)(2.0,0.46)(2.2,0.66)(2.4,0.84)(2.6,0.94)(2.8,0.96)};
  \draw[blue!70,thick,smooth] plot coordinates
    {(0,0.73)(0.2,0.78)(0.4,0.97)(0.6,0.96)(0.8,0.98)(1.0,0.72)(1.2,0.58)
     (1.4,0.32)(1.6,0.29)(1.8,0.26)(2.0,0.50)(2.2,0.69)(2.4,0.81)(2.6,0.97)(2.8,0.93)};
  \node[above,font=\footnotesize] at (1.4,1.3) {低噪声（\(t\) 小）};
\end{scope}
% --- 中噪声 ---
\begin{scope}[shift={(3.8,2.3)}]
  \draw[black!12] (0,0) rectangle (2.8,1.2);
  \draw[black!30,thin,smooth] plot coordinates
    {(0,0.70)(0.2,0.82)(0.4,0.94)(0.6,1.00)(0.8,0.94)(1.0,0.76)(1.2,0.54)
     (1.4,0.36)(1.6,0.26)(1.8,0.30)(2.0,0.46)(2.2,0.66)(2.4,0.84)(2.6,0.94)(2.8,0.96)};
  \draw[violet!70,thick,smooth] plot coordinates
    {(0,0.88)(0.2,0.65)(0.4,1.10)(0.6,0.85)(0.8,1.05)(1.0,0.94)(1.2,0.32)
     (1.4,0.46)(1.6,0.10)(1.8,0.44)(2.0,0.26)(2.2,0.82)(2.4,0.70)(2.6,1.10)(2.8,0.80)};
  \node[above,font=\footnotesize] at (1.4,1.3) {中噪声};
\end{scope}
% --- 高噪声 ---
\begin{scope}[shift={(7.6,2.3)}]
  \draw[black!12] (0,0) rectangle (2.8,1.2);
  \draw[black!30,thin,smooth] plot coordinates
    {(0,0.70)(0.2,0.82)(0.4,0.94)(0.6,1.00)(0.8,0.94)(1.0,0.76)(1.2,0.54)
     (1.4,0.36)(1.6,0.26)(1.8,0.30)(2.0,0.46)(2.2,0.66)(2.4,0.84)(2.6,0.94)(2.8,0.96)};
  \draw[red!70,thick,smooth] plot coordinates
    {(0,1.02)(0.2,0.25)(0.4,0.90)(0.6,0.50)(0.8,1.08)(1.0,0.16)(1.2,0.80)
     (1.4,0.30)(1.6,0.65)(1.8,1.05)(2.0,0.20)(2.2,0.90)(2.4,0.40)(2.6,1.00)(2.8,0.55)};
  \node[above,font=\footnotesize] at (1.4,1.3) {高噪声（\(t\) 大）};
\end{scope}
% ===== 汇聚箭头 =====
\draw[->,blue!60,thick]   (1.4,2.2) -- (4.5,1.15);
\draw[->,violet!60,thick] (5.2,2.2) -- (5.2,1.15);
\draw[->,red!60,thick]    (9.0,2.2) -- (5.9,1.15);
% ===== 网络方块 =====
\fill[black!6] (3.9,0.45) rectangle (6.5,1.1);
\draw[black!55] (3.9,0.45) rectangle (6.5,1.1);
\node at (5.2,0.775) {\(\varepsilon_\theta(\,\cdot\,,t)\)};
% t 作为输入
\draw[->,black!55] (2.75,0.775) -- (3.85,0.775);
\node[left,font=\footnotesize] at (2.7,0.775) {噪声水平 \(t\) 也是输入};
% ===== 三条任务说明 =====
\fill[blue!70]   (0.15,-0.05) rectangle (0.35,0.15);
\node[right,font=\footnotesize] at (0.5,0.05) {低噪声任务：结构还在，要补的是纹理、边缘锐度、局部一致性——\emph{细节尺度}};
\fill[violet!70] (0.15,-0.55) rectangle (0.35,-0.35);
\node[right,font=\footnotesize] at (0.5,-0.45) {中噪声任务：轮廓半毁，要恢复的是部件位置与相互关系——\emph{中层结构}};
\fill[red!70]    (0.15,-1.05) rectangle (0.35,-0.85);
\node[right,font=\footnotesize] at (0.5,-0.95) {高噪声任务：只剩大致明暗，能定的只有构图与整体色调——\emph{全局尺度}};
\end{tikzpicture}

\emph{三幅图里灰色的干净信号是同一条，只是被不同强度的噪声盖住。三个箭头指向同一个网络方块：参数共享，噪声水平 \(t\) 作为输入进去。要记住的是下面三行——同一组权重被要求在三种尺度上分别做不同的事，网络必须依靠 \(t\) 把自己切换到对应的模式。}

\end{center}

这解释了一个在实践中反复被观察到的现象：采样早期（高噪声）决定构图，后期（低噪声）决定细节。原因就在图里——早期调用的是网络的全局那一档能力，那时它能给出的信息本来就只有大尺度结构；等噪声降下去，构图已经定死了，剩下的容量全用在填细节上。改提示词只影响早期几步，改随机种子会换掉整张图的布局，都是同一件事的表现。

这个多任务的性质也决定了架构上的一些选择。\(t\) 通常先编码成一组正弦特征，再注入到网络的每一层，而不是拼在输入通道上——它需要在所有分辨率上都能改变网络的行为，只在入口喂一次做不到这一点。U-Net 那种跨尺度的骨架在这里合适，也是因为不同噪声水平的任务天然落在不同的分辨率层级上。这些都是经验上收敛下来的做法，没有哪条定理说必须如此。

至此模型学什么已经清楚了：一个跨所有噪声水平的去噪回归器，等价地，一族得分估计。但有了得分场不等于有了样本。从纯噪声出发，沿着这个场往回走，每一步走多远、走多少步、要不要在中途继续注入随机性——这些问题里藏着扩散模型真正的成本。下一节处理它们。
